% from ega4-2-fr.tex
\oldpage[IV-2]{179}
\begin{proposition}[6.15.6]
\label{IV.2.s6.15.6-fr}
Soient $k$ un corps, $X$ un $k$-préschéma. Si $X$ est normal, alors, pour
toute extension $k'$ de $k$, $X'=X\otimes_k k'$ est géométriquement
unibranche.
\end{proposition}

\begin{proof}
On sait que $k'$ est extension algébrique d'une extension pure $k_0$ de
$k$, et si $k''$ est la plus grande extension séparable de $k_0$ contenue
dans $k'$, $k'$ est une extension radicielle de $k''$ et $k''$ extension
séparable de $k$. On sait (6.14.2) que $X''=X\otimes_k k''$ est normal;
comme $X\otimes_k k'=X''\otimes_{k''}k'$, on voit qu'on peut se borner au
cas où l'extension $k'$ de $k$ est radicielle. En outre (I, 3.6.5), on peut
supposer que $X=\operatorname{Spec}(A)$, où $A$ est un anneau local intègre
et intégralement clos (puisque $X$ est normal). Le morphisme projection
$f:X'\to X$ est un homéomorphisme, puisque
$\operatorname{Spec}(k')\to\operatorname{Spec}(k)$ est un homéomorphisme
universel (2.4.5); comme $X'=\operatorname{Spec}(A')$ où
$A'=A\otimes_k k'$, on voit donc que $A'$ est un anneau local dont le
nilradical $\mathfrak N'$ est l'unique idéal premier minimal, d'où
$X'_{\mathrm{red}}=\operatorname{Spec}(A'_0)$, où
$A'_0=A'/\mathfrak N'$ est un anneau local intègre; en outre, si $K$ est
le corps des fractions de $A$, le corps des fractions $K'_0$ de $A'_0$ est
radiciel sur $K$, puisque le morphisme $f$ est radiciel. Comme $A'_0$ est
entier sur $A$, sa clôture intégrale $B$ est aussi la fermeture intégrale
de $A$ dans $K'_0$. Mais comme $A$ est intégralement clos, on sait
(Bourbaki, Alg. comm., chap. V, § 2, n$^o$ 3, lemme 3) que $B$ est
l'ensemble des $x\in K'_0$ dont une puissance $p^m$-ème (pour un $m$ assez
grand) appartient à $A$ ($p$ étant l'exposant caractéristique de $K$);
en outre il n'existe qu'un seul idéal premier de $B$ au-dessus de tout idéal
premier de $A$; en particulier $B$ est un anneau local et son corps résiduel
est une extension radicielle de celui de $A$, et a fortiori de celui de
$A'_0$, ce qui prouve que $A'_0$ est géométriquement unibranche, et par
suite il en est de même de $X'$.
\end{proof}

\begin{proposition}[6.15.7]
\EGAReaderTarget{ega4.c04.prop.6-15-7}{IV.2.s6.15.7-fr}
Soient $k$ un corps, $X$ un $k$-préschéma, $k'$ une extension de $k$,
$X'=X\otimes_k k'$. Soient $x'$ un point de $X'$, $x$ sa projection dans
$X$. Pour que $X$ soit géométriquement unibranche au point $x$, il faut et
il suffit que $X'$ soit géométriquement unibranche au point $x'$. Pour que
$X$ soit géométriquement unibranche, il faut et il suffit que $X'$ le soit.
\end{proposition}

\begin{proof}
La seconde assertion résulte de la première et du fait que la projection
$f:X'\to X$ est un morphisme surjectif. Pour démontrer la première
assertion, on peut, en vertu de (I, 5.1.8), se borner au cas où $X$ est
réduit, et en vertu de (I, 3.6.5), au cas
